\section{Overview and System Contract}
\label{sec:overview}

\subsection{Three optional mechanisms over a common file interface}

The common path provides file operations, metadata resolution, and object
access. Three choices refine that path at different boundaries
(Table~\ref{tab:mechanism-choices}). View composition determines which state
a task sees. Metadata representation determines how immutable records are
located and fetched. Object granularity determines how payloads are formed
and how much work their access requires. No choice semantically requires
the other two.

\begin{table*}[t]
\centering
\small
\begin{tabularx}{\textwidth}{@{}p{3.0cm}YY@{}}
\toprule
Choice & Mechanism selected & Alternative retained\\
\midrule
View composition & Delta-based workspace: shared baseline plus private state & Ordinary volume with its own namespace and mappings\\
Immutable metadata & Packed Metadata: indexed records in bounded object-resident pages & Native catalog or KV representation\\
Object granularity & Configured client block size $g$; optional compatible backend $B/q$ choices & Fixed client geometry and backend default layout\\
\bottomrule
\end{tabularx}
\caption{Independent design decisions. A deployment selects mechanisms for
their own costs and contracts. Current mount paths support specific
combinations; this table is not a claim of eight interchangeable runtime modes.}
\label{tab:mechanism-choices}
\end{table*}

For example, a task requiring independent edits can use a delta-based
workspace over a native immutable layer. A read-only environment can use
a packed snapshot directly. An ordinary writable volume can configure its
object block size without either facility. When repeated tasks share a stable
environment, they can instead combine workspaces and a packed baseline.
Optional backend layout selection is another data-path decision in each
eligible uploader, rather than a condition for any of these view types.

\subsection{A representation-independent workspace contract}

Let $S$ identify a shared immutable baseline and $\Delta_i$ the private
writable state of workspace $i$. Its view is
\begin{equation}
  W_i = \operatorname{Resolve}(S,\Delta_i).
\end{equation}
The baseline exposes name, attribute, and data-location queries. It may bind
a sealed native layer or a packed manifest; $S$ does not mean that every
workspace has a packed lower. Clients pin an identity instead of following
a mutable ``latest'' pointer on each lookup. A selected subtree can form a
baseline without freezing an entire source volume.

The upper represents private names, attributes, deletion state, and data
mappings. Newly written payload bytes reside in separately referenced objects.
Resolution checks private state first and inherits from $S$ only when the
upper permits it. A deletion marker suppresses a lower entry, while absence
of an upper record permits inheritance. This contract leaves the lower's
physical metadata representation open.

The isolation property is view-local mutation: an operation in $W_i$ leaves
$S$ and another workspace's private state unchanged. This property does not
provide process isolation, network restrictions, or automatic merging of
independent edits. Those functions require their own runtime or application
protocols. Concurrent operations within a workspace must separately preserve
its namespace and data invariants.

\subsection{Composition through explicit interfaces}

Figure~\ref{fig:overview} illustrates one possible composition using an
explicit upper data map, a packed lower, and optional backend layout hints.
The figure expands the interfaces between mechanisms; it is not the minimum
configuration. For example, two workspaces can share a fetched baseline page,
while one records a private overwrite and the other continues to read $S$.
The chosen lower reader resolves inherited data locators to payload ranges.
The backend serves those objects using their persisted storage layout.

\begin{figure*}[!t]
\centering
\begin{tikzpicture}[
  box/.style={draw,rounded corners=2pt,align=center,text width=4.5cm,minimum height=0.85cm,font=\small,fill=black!3},
  flow/.style={-{Latex[length=2mm]},semithick},
  todo/.style={flow,dashed},
  lab/.style={font=\footnotesize,fill=white,inner sep=1.5pt}
]
\node[box] (view) at (0,0) {Agent tools and FUSE view\\$W_i=\operatorname{Resolve}(S,\Delta_i)$};
\node[box] (upper) at (-5.3,-1.55) {Private upper $\Delta_i$\\names, attributes, data maps};
\node[box] (lower) at (0,-1.55) {Packed lower reader\\name ranges and paged indexes};
\node[box] (producer) at (5.3,-1.55) {Snapshot producer\\selected immutable subtrees};
\node[box] (objects) at (0,-3.2) {Shared snapshot objects\\manifest, indexes, GroupMeta, frames};
\node[box] (upload) at (-5.3,-3.2) {Private data uploads\\finalized object length $L_o$};
\node[box] (backend) at (5.3,-3.2) {RustFS object layout\\persisted coding $B$ and read frame $q$};
\draw[flow] (view) -| (upper);
\draw[flow] (view) -- (lower);
\draw[flow] (upper) -- (upload);
\draw[flow] (lower) -- (objects);
\draw[flow] (producer) -- (objects);
\draw[flow] (objects) -- (backend);
\draw[todo] (producer) -- node[lab,right]{layout hint} (backend);
\draw[todo] (upload.south) -- ++(0,-0.65) -| node[lab,pos=0.25,below]{upper upload hint: integration gate} (backend.south);
\end{tikzpicture}

\caption{An illustrative combination of the optional mechanisms.
(a) An Entity-style upper resolves private state and explicit Data, LowerRange,
and Hole intervals. (b) A packed lower resolves a representative name lookup
through an index and GroupMeta page to a payload descriptor. A native lower
can implement the baseline contract instead; Packed also supports standalone
read-only use. (c) Finalized objects may request a compatible backend layout;
omitted hints retain default selection. Client object formation with $g$
precedes this panel and is detailed in Section~\ref{sec:layout}.
Solid arrows show access dependencies within this combination; dashed hint
wiring remains incomplete in the illustrated Entity/producer paths. The 2+2
shards are schematic. This is a design composition, not an end-to-end result.}
\label{fig:overview}
\end{figure*}

Composition preserves three boundaries. First, transient dirty buffers may
provide local read-after-write behavior before objects are remotely durable;
they are distinct from the workspace's private namespace state. Second,
publication makes a validated immutable version visible only after its
required objects satisfy the selected persistence contract. Reader ownership
must keep referenced state alive through reclamation. Third, logical object
identity and file contents remain independent of the backend's coding layout.
Changing a backend setting does not change an already published view.

A recoverable file version can include explicit application checkpoints,
but does not reconstruct process memory, open handles, or execution state
by itself. Section~\ref{sec:implementation} distinguishes current mount
paths and recovery support. The following sections develop the principles
of each mechanism before those implementation-specific limits.
